\documentclass[10pt,a4paper]{amsart}
\usepackage[margin=19mm]{geometry}
\usepackage{amsmath,amssymb,mathtools}
\usepackage[scr=rsfs,cal=euler]{mathalfa}
\usepackage{xfrac}
\usepackage[unicode,hidelinks,bookmarks=false]{hyperref}
\newcommand{\ie}{i.e.\ }
\newcommand{\cf}{cf.\ }
\newcommand{\resp}{respectively\ }
\newcommand{\Subsection}{\subsection}
\providecommand{\nobreakdash}{\nobreak\hbox{-}}
\setlength{\emergencystretch}{2em}
\multlinegap0pt
\usepackage{mathtools}
\usepackage{amssymb}
\usepackage{xcolor}
\usepackage{dsfont}
\usepackage{tikz}
\usepackage{bm}
\usepackage[shortlabels]{enumitem}
\newtheorem{theorem}{Theorem}
\newtheorem{corollary}{Corollary}
\def\EE{\mathbb{E}}
\def\PP{\mathbb{P}}
\def\RR{\mathbb{R}}
\newcommand{\dint}{\textup{d}}
\newcommand{\vol}{\mathop{\mathrm{vol}}\nolimits}
\title{\bfseries Visibility in the Boolean Model \\ on Harmonic Manifolds}
\author{Enkelejd Hashorva, Christoph Th\"ale}
\date{Source: arXiv:2605.23817v1 (CC BY 4.0)}
\begin{document}
\maketitle

\noindent\emph{Source and licence.} \bfseries Visibility in the Boolean Model \\ on Harmonic Manifolds. Version arXiv:2605.23817v1 (CC BY 4.0), distributed by arXiv under the Creative Commons Attribution 4.0 International licence, http://creativecommons.org/licenses/by/4.0/, as recorded in the arXiv licence metadata for this exact version and shown on the abstract page https://arxiv.org/abs/2605.23817v1. Full licence notice: \texttt{licenses/arxiv-2605.23817-CC-BY-4.0.txt}.\par\smallskip

\noindent\textit{Editorial scope.} Counted unit: the article's corollary cor (source0.tex lines 610--621). The article's complete original proof of it is delivered with it (source0.tex lines 622--664, one \texttt{proof} environment of 43 lines). No mathematics is written, repaired or re-derived: the statement and the proof are the article's own lines, in the article's own notation, and no retained line is substituted or re-flowed.  Retained uncounted as the source-local context the proof needs: lines 492--497; lines 549--561.  Nothing was omitted from any retained span, so the package carries no omission record.  The retained lines cite 2 works, whose entries are transcribed verbatim from the article's own bibliography source in the e-print and delivered with short editorial labels recorded in the item's reference mapping.  No retained line contains a figure environment, an image-inclusion command or drawing code, so no image is delivered and the package declares no asset dependency.  Editorial formatting only, in addition: an AMS \texttt{amsart} preamble that restates the article's own declarations and macros used by the retained lines, namely the environments \texttt{theorem}, \texttt{corollary} on separate counters, together with the standard packages those lines need.  The retained environments print in this document as Theorem 1, Corollary 1 in order of appearance, whereas the article prints the same items with the numbering of its own full document.  The complete native source of the article as distributed in the arXiv e-print for this exact version is bundled unchanged as \texttt{sources/201237/source0.tex}.
\begin{equation}\label{eq:arho}
	\vol_M(T_\rho(\gamma_u[0,r]))
	=
	\vol_M(B(o,\rho))+a_\rho r,
	\qquad r\geq 0.
\end{equation}

\begin{theorem}\label{thm:exp_visibility_harmonic}
	Let \(M\) be a simply connected non-compact homogeneous harmonic manifold.
	Consider the Boolean model on \(M\) with intensity \(\lambda>0\) and fixed
	ball grains of radius \(\rho>0\). Then there exists a constant
	\(a_\rho\in(0,\infty)\), depending only on \(M\) and \(\rho\), such that
	for every \(o\in M\), every unit tangent vector \(u\in T_oM\), and every
	\(r\geq 0\),
	\[
	\PP\bigl(s(u)>r\,\big|\,o\notin Z\bigr)
	=
	e^{-\lambda a_\rho r}.
	\]
\end{theorem}

\begin{corollary}\label{cor}
	Let \(M\) be a simply connected non-compact homogeneous harmonic manifold
	with volume entropy \(h\geq 0\). Consider the Boolean model on \(M\) with
	intensity \(\lambda>0\) and fixed ball grains of radius \(\rho>0\), and let
	\(a_\rho\) be the tube coefficient from \eqref{eq:arho}. If \(h=0\), then $\EE[\vol_M(V_o)\mid o\notin Z]<\infty$ for every $\lambda>0$.
	If \(h>0\), then
	\[
	\EE[\vol_M(V_o)\mid o\notin Z]<\infty
	\qquad\Longleftrightarrow\qquad
	\lambda>\lambda_c(M,\rho):={h\over a_\rho}.
	\]
\end{corollary}

\begin{proof}
	By Fubini's theorem,
	\[
	\EE[\vol_M(V_o)\mid o\notin Z]
	=
	\int_M
	\PP(x\in V_o\mid o\notin Z)\,\vol_M(\dint x).
	\]
	Since the spaces
	under consideration are simply connected homogeneous harmonic manifolds, we may
	use global geodesic polar coordinates around \(o\). Then
	\begin{equation}\label{eq:intThreshold}
	\EE[\vol_M(V_o)\mid o\notin Z]
	=
	\int_0^\infty
	e^{-\lambda a_\rho r}S(r)\,\dint r,
	\end{equation}
	where \(S(r)\) denotes the surface area of the
	geodesic sphere of radius \(r\) around \(o\). Indeed, by
	Theorem~\ref{thm:exp_visibility_harmonic}, the conditional probability that
	the geodesic segment of length \(r\) from \(o\) is visible is
	\(e^{-\lambda a_\rho r}\), independently of the direction.
	
	If \(h=0\), then \(M\) is flat by Heber's classification in \cite{HeberHarmonicHomogeneous}, and hence
	\(M\) is isometric to \(\RR^n\). In this case $S(r)=n\kappa_n r^{n-1}$,
	and therefore $\int_0^\infty e^{-\lambda a_\rho r}S(r)\,\dint r<\infty$ for every \(\lambda>0\).
	
	It remains to consider the case \(h>0\). For simply connected non-compact
	homogeneous harmonic manifolds with positive entropy, the geodesic sphere
	volume has pure exponential growth with exponent \(h\), i.e., there exist
	constants \(c_1,c_2>0\) such that
	\[
	c_1e^{hr}\leq S(r)\leq c_2e^{hr},
	\qquad r\geq 1.
	\]
	This follows from the known volume density estimates for these spaces, see
	for instance the discussion before Definition~2.2 and Corollary~5.3 in
	\cite{Knieper12}. Hence, the tail behaviour of the integral in \eqref{eq:intThreshold} is the same as that of
	\[
	\int_{1}^\infty e^{-(\lambda a_\rho-h)r}\,\dint r.
	\]
	This integral is finite if and only if \(\lambda a_\rho>h\). This proves the result.
\end{proof}

\begin{thebibliography}{99}
\bibitem[HeberH]{HeberHarmonicHomogeneous} Heber, J.: On harmonic and asymptotically harmonic homogeneous spaces. \textit{Geometric and Functional Analysis} \textbf{16}, 869--890 (2006).
\bibitem[Kniepe]{Knieper12} Knieper, G.: New results on noncompact harmonic manifolds. \textit{Commentarii Mathematici Helvetici} \textbf{87}, 669--703 (2012).
\end{thebibliography}
\end{document}
